\section{Riemann 流形}\label{sec:riemannian-metrics}

为定义 Riemann 流形, 必须先更仔细地考察 2-张量.

\begin{RiemannDefinition}\label{def:symmetric-nondegenerate-positive-tensor}
张量 $g\in\mathcal{T}^2(T_p^*M)$ 称为
\begin{enumerate}[label=(\roman*),leftmargin=*]
\item \textbf{对称的}, 若对所有 $v,w\in T_pM$, 有 $g(v,w)=g(w,v)$;
\item \textbf{非退化的}, 若对所有 $w\in T_pM$ 有 $g(v,w)=0$ 蕴含 $v=0$;
\item \textbf{正定的}, 若对所有 $v\in T_pM\setminus\{0\}$, 有 $g(v,v)>0$.
\end{enumerate}
\end{RiemannDefinition}

若对所有 $p\in M$, 都有 $g_p$ 对称, 非退化或正定, 则称协变 2-张量场 $g$ 对称, 非退化或正定. 若 $x:V\to\mathbb{R}^n$ 是一个局部坐标图, 则在 $V$ 上有
\[
  g=\sum_{i,j=1}^n g_{ij}\,\mathrm{d}x^i\otimes\mathrm{d}x^j
\]
其中
\[
  g_{ij}=g\left(\frac{\partial}{\partial x^i},\frac{\partial}{\partial x^j}\right).
\]
容易看出, $g$ 对称, 非退化或正定, 当且仅当矩阵 $(g_{ij})$ 具有相应的性质 [见 Exercise~\ref{ex:riemannian-manifolds}(\ref{ex:matrix-criteria-for-metric-tensor})].

\begin{RiemannDefinition}\label{def:riemannian-metric-and-manifold}
光滑流形 $M$ 上的一个 \textbf{Riemann 度量}是对称正定的光滑协变 2-张量场 $g$. 配备了 Riemann 度量 $g$ 的光滑流形 $M$ 称为 \textbf{Riemann 流形}, 记作 $(M,g)$.
\end{RiemannDefinition}

因此, Riemann 度量以光滑的方式给每个切空间指定了一个内积. 通常记作
\[
  g_p(v,w)=\langle v,w\rangle_p.
\]

\begin{RiemannExample}[{\unboldmath Euclidean $n$ 维空间}]\label{ex:euclidean-metric}
显然, $M=\mathbb{R}^n$ 和
\[
  g=\sum_{i=1}^n\mathrm{d}x^i\otimes\mathrm{d}x^i
\]
定义了一个 Riemann 流形.
\end{RiemannExample}

\begin{RiemannProposition}\label{prop:induced-metric-of-immersion}
设 $(N,g)$ 是 Riemann 流形, $f:M\to N$ 是浸入. 则 $f^*g$ 是 $M$ 上的一个 Riemann 度量 (称为\textbf{诱导度量}).
\end{RiemannProposition}

\begin{RiemannProof}
只需证明 $f^*g$ 对称且正定. 设 $p\in M$ 且 $v,w\in T_pM$. 由于 $g$ 对称,
\begin{align*}
  (f^*g)_p(v,w)&=g_{f(p)}((\mathrm{d}f)_pv,(\mathrm{d}f)_pw)\\
  &=g_{f(p)}((\mathrm{d}f)_pw,(\mathrm{d}f)_pv)=(f^*g)_p(w,v).
\end{align*}
另一方面, 显然 $(f^*g)_p(v,v)\geq0$, 且
\[
  (f^*g)_p(v,v)=0\Rightarrow g_{f(p)}((\mathrm{d}f)_pv,(\mathrm{d}f)_pv)=0
  \Rightarrow(\mathrm{d}f)_pv=0\Rightarrow v=0
\]
(因为 $(\mathrm{d}f)_p$ 是单射).
\end{RiemannProof}

特别地, Riemann 流形 $(N,g)$ 的任意子流形 $M$ 本身就是 Riemann 流形. 注意, 在这种情形下, 每个点 $p\in M$ 处的诱导度量恰为 $g_p$ 在 $T_pM\subset T_pN$ 上的限制. 由于 $\mathbb{R}^n$ 是 Riemann 流形 (参见 Example~\ref{ex:euclidean-metric}), $\mathbb{R}^n$ 的任意子流形都是 Riemann 流形. 于是 Whitney 定理蕴含任意流形都容许一个 Riemann 度量.

John Nash 在 1954 年证明 \RGcite{Nas56}: 任意紧的 $n$ 维 Riemann 流形都可等距嵌入 $\mathbb{R}^N$, 其中 $N=\frac{n(3n+11)}2$ (也就是说, 嵌入使其度量由 $\mathbb{R}^N$ 的 Euclidean 度量诱导). 后来 Gromov \RGcite{GR70} 证明可以取 $N=\frac{(n+2)(n+3)}2$. 注意, 当 $n=2$ 时, Nash 的结果给出任意紧曲面在 $\mathbb{R}^{17}$ 中的等距嵌入, 而 Gromov 的结果给出在 $\mathbb{R}^{10}$ 中的等距嵌入. 事实上, Gromov 进一步证明任意曲面都能等距嵌入 $\mathbb{R}^5$. 这个结果不能改进, 因为带标准度量的实射影平面 [见 Exercise~\ref{ex:riemannian-manifolds}(\ref{ex:quotient-riemannian-structures})] 不能等距嵌入 $\mathbb{R}^4$.

\begin{RiemannExample}\label{ex:standard-metric-of-sphere}
\[
  S^n=\{x\in\mathbb{R}^{n+1}\mid\lVert x\rVert=1\}
\]
上的\textbf{标准度量}是由 $\mathbb{R}^{n+1}$ 上的 Euclidean 度量在 $S^n$ 上诱导的度量. 例如, 开集
\[
  U=\{x\in S^n\mid x^{n+1}>0\}
\]
的一个参数化为
\[
  \varphi(x^1,\ldots,x^n)=\left(x^1,\ldots,x^n,\sqrt{1-(x^1)^2-\cdots-(x^n)^2}\right),
\]
相应的度量张量系数为
\[
  g_{ij}=\left\langle\frac{\partial\varphi}{\partial x^i},\frac{\partial\varphi}{\partial x^j}\right\rangle
  =\delta_{ij}+\frac{x^ix^j}{1-(x^1)^2-\cdots-(x^n)^2}.
\]
\end{RiemannExample}

若两个 Riemann 流形等距, 则将它们视为相同.

\begin{RiemannDefinition}\label{def:isometry-and-local-isometry}
设 $(M,g)$ 和 $(N,h)$ 是 Riemann 流形. 若微分同胚 $f:M\to N$ 满足 $f^*h=g$, 就称它为\textbf{等距映射}. 类似地, 若局部微分同胚 $f:M\to N$ 满足 $f^*h=g$, 就称它为\textbf{局部等距映射}.
\end{RiemannDefinition}

借助 Riemann 度量, 可以计算任意向量 $v\in TM$ 的\textbf{长度} $\lVert v\rVert=\langle v,v\rangle^{\frac12}$ (也可以计算两个基点相同的向量之间的\phantomsection\label{idx:angle-between-tangent-vectors}\textbf{夹角}). 因而可以度量曲线的长度.

\begin{RiemannDefinition}\label{def:length-of-curve}
若 $(M,\langle\cdot,\cdot\rangle)$ 是 Riemann 流形, $c:[a,b]\to M$ 是可微曲线, 则 $c$ 的\textbf{长度}为
\[
  l(c)=\int_a^b\lVert\dot c(t)\rVert\,\mathrm{d}t.
\]
\end{RiemannDefinition}

一段曲线的长度与参数化无关 [见 Exercise~\ref{ex:riemannian-manifolds}(\ref{ex:length-independent-of-parameterization})].

回顾, 若 $M$ 是 $n$ 维可定向流形, 则它具有体积元, 即对所有 $p\in M$ 都有 $\omega_p\ne0$ 的微分形式 $\omega\in\Omega^n(M)$. 显然, 体积元与无零点的可微函数 $f\in C^\infty(M)$ 一样多.

\begin{RiemannDefinition}\label{def:riemannian-volume-element}
若 $(M,g)$ 是可定向 Riemann 流形, 对每个 $p\in M$ 及 $T_pM$ 的任意标准正交基 $\{v_1,\ldots,v_n\}$, 若
\[
  \omega_p(v_1,\ldots,v_n)=\pm1
\]
则称 $\omega\in\Omega^n(M)$ 为 \textbf{Riemann 体积元}.
\end{RiemannDefinition}

注意, 若 $M$ 连通, 则恰有两个 Riemann 体积元 (每种定向的选取各对应一个). 此外, 若 $\omega$ 是 Riemann 体积元, 且 $x:V\to\mathbb{R}$ 是与 $\omega$ 诱导的定向相容的坐标图, 则
\[
  \omega=f\,\mathrm{d}x^1\wedge\ldots\wedge\mathrm{d}x^n
\]
其中有正函数
\[
  f=\omega\left(\frac{\partial}{\partial x^1},\ldots,\frac{\partial}{\partial x^n}\right).
\]
若 $S$ 是以 $\frac{\partial}{\partial x^1},\ldots,\frac{\partial}{\partial x^n}$ 在某组具有相同定向的标准正交基下的分量为列的矩阵, 则有
\[
  f=\det S=\bigl(\det(S^2)\bigr)^{\frac12}
  =\bigl(\det(S^tS)\bigr)^{\frac12}=\bigl(\det(g_{ij})\bigr)^{\frac12}
\]
因为 $S^tS$ 的第 $(i,j)$ 项显然是内积 $g\bigl(\frac{\partial}{\partial x^i},\frac{\partial}{\partial x^j}\bigr)=g_{ij}$.

$M$ 上的 Riemann 度量 $\langle\cdot,\cdot\rangle$ 对所有 $p\in M$ 确定一个线性同构 $\Phi_g:T_pM\to T_p^*M$, 由 $\Phi_g(v)(w)=\langle v,w\rangle$ 对所有 $v,w\in T_pM$ 定义. 它延伸为 $\mathfrak{X}(M)$ 和 $\Omega^1(M)$ 之间的同构. 特别地, 有如下定义.

\begin{RiemannDefinition}\label{def:gradient-of-function}
设 $(M,g)$ 是 Riemann 流形, $f:M\to\mathbb{R}$ 是光滑函数. $f$ 的\textbf{梯度}是通过由 $g$ 确定的同构与 1-形式 $\mathrm{d}f$ 相对应的向量场 $\operatorname{grad}f$.
\end{RiemannDefinition}

\begin{RiemannExercise}\label{ex:riemannian-manifolds}
\begin{enumerate}[label=(\arabic*),ref=\arabic*,leftmargin=*]
\item\label{ex:matrix-criteria-for-metric-tensor} 设 $g=\sum_{i,j=1}^n g_{ij}\,\mathrm{d}x^i\otimes\mathrm{d}x^j\in\mathcal{T}^2(T_p^*M)$. 证明:
\begin{enumerate}[label=(\alph*),leftmargin=*]
\item $g$ 对称当且仅当 $g_{ij}=g_{ji}$ ($i,j=1,\ldots,n$);
\item $g$ 非退化当且仅当 $\det(g_{ij})\ne0$;
\item $g$ 正定当且仅当 $(g_{ij})$ 是正定矩阵;
\item 若 $g$ 非退化, 则由 $\Phi_g(v)(w)=g(v,w)$ 对所有 $v,w\in T_pM$ 给出的映射 $\Phi_g:T_pM\to T_p^*M$ 是线性同构;
\item 若 $g$ 正定, 则 $g$ 非退化.
\end{enumerate}
\item\label{ex:riemannian-structure-by-partition-of-unity} 不借助 Whitney 定理, 证明任意可微流形都容许 Riemann 结构.
{\small (\textbf{提示:} 使用单位分解).}
\item\label{ex:quotient-riemannian-structures}
\begin{enumerate}[label=(\alph*),leftmargin=*]
\item 设 $(M,g)$ 是 Riemann 流形, $G$ 是通过等距映射自由且固有地作用在 $M$ 上的离散 Lie 群. 证明 $M/G$ 具有自然的 Riemann 结构 (称为\textbf{商结构}).
\item 如何定义 $n$ 维环面 $T^n=\mathbb{R}^n/\mathbb{Z}^n$ 上的\textbf{平直方形度量}?
\item 如何定义 $n$ 维实射影空间 $\mathbb{R}P^n=S^n/\mathbb{Z}_2$ 上的\textbf{标准度量}?
\end{enumerate}
\item\label{ex:left-invariant-metric} 回顾, 给定 Lie 群 $G$ 和 $x\in G$, 由 $L_x(y)=xy$ 对所有 $y\in G$ 给出的微分同胚 $L_x:G\to G$ 是由 $x$ 确定的左平移. 若对所有 $x\in G$, $L_x$ 都是等距映射, 则称 $G$ 上的 Riemann 度量 $g$ \textbf{左不变}. 证明:
\begin{enumerate}[label=(\alph*),leftmargin=*]
\item $g(\cdot,\cdot)\equiv\langle\cdot,\cdot\rangle$ 左不变当且仅当
\[
  \langle v,w\rangle_x=\bigl\langle(\mathrm{d}L_{x^{-1}})_xv,(\mathrm{d}L_{x^{-1}})_xw\bigr\rangle_e
\]
对所有 $x\in G$ 和 $v,w\in T_xG$ 成立, 其中 $e\in G$ 为单位元, $\langle\cdot,\cdot\rangle_e$ 是 Lie 代数 $\mathfrak{g}=T_eG$ 上的内积;
\item $S^3\cong SU(2)$ 上的标准度量左不变;
\item $\mathcal{M}_{n\times n}\cong\mathbb{R}^{n^2}$ 的 Euclidean 度量在 $O(n)$ 上诱导的度量左不变.
\end{enumerate}
\item\label{ex:length-independent-of-parameterization} 若存在光滑双射 $f:[\alpha,\beta]\to[a,b]$ (即重参数化), 使得 $\gamma=c\circ f$, 就说可微曲线 $\gamma:[\alpha,\beta]\to M$ 由曲线 $c:[a,b]\to M$ \textbf{重参数化}得到. 证明若 $\gamma$ 由 $c$ 重参数化得到, 则 $l(\gamma)=l(c)$.
\item\label{ex:gradient-normal-to-regular-level-set} 设 $(M,g)$ 是 Riemann 流形, $f\in C^\infty(M)$. 证明若 $a\in\mathbb{R}$ 是 $f$ 的正则值, 则 $\operatorname{grad}(f)$ 与子流形 $f^{-1}(a)$ 正交.
\end{enumerate}
\end{RiemannExercise}
